
\subsubsection*{2.6.2}
\phantomsection
\label{sga3:VIIB:2.6.2}

Conversely, every Lie algebra \(\mathbf L\) over \(\mathbf O_k\)
defines a \(k\)-group functor. Indeed, let \(U(\mathbf L)\) denote
the functor assigning to every object \(C\) of
\(\operatorname{Alf}/k\) the enveloping algebra
\[
  U\bigl(\mathbf L(C)\bigr)
\]
of the \(C\)-Lie algebra \(\mathbf L(C)\). By
VII\textsubscript{A}, \SGARefLink{sga3:VIIA:3.2.2}{3.2.2}, \(U(\mathbf L)\) is a bialgebra over
\(\mathbf O_k\), and hence, by \SGARefLink{sga3:VIIA:2.2}{2.2}, determines the \(k\)-group
functor
\[
  \operatorname{Spf}^{*}U(\mathbf L),
\]
which from now on we denote by \(\mathscr G(\mathbf L)\). Thus
\(\mathscr G(\mathbf L)(C)\) is the group of elements
\(z\in U(\mathbf L(C))\) having augmentation \(1\) and satisfying
\[
  \Delta_{U(\mathbf L(C))}(z)=z\otimes z.
\]

When \(\mathbf L\) is flat over \(\mathbf O_k\), we have the
following proposition.

\medskip
\noindent\textbf{Proposition.}\footnote{%
\label{sga3:VIIB:nde:110}The editors set off parts \textup{\SGARefLink{sga3:vii:paragraph:5-1:equation:i}{(i)}} and
\textup{\SGARefLink{sga3:vii:paragraph:5-1:equation:ii}{(ii)}} and added part \textup{(iii)}, which will be used in
\SGARefLink{sga3:VIIB:2.6.3:prop}{2.6.3} and \SGARefLink{sga3:VIIB:3.3.2}{3.3.2}.}
\phantomsection
\label{sga3:VIIB:2.6.2:prop}
{\itshape Let \(\mathbf L\) be a flat Lie algebra over
\(\mathbf O_k\).
\begin{enumerate}
\item[\textup{\SGAInlineRefTarget{sga3:VIIB:2.6.2:prop:item:i}(i)}]
\(\mathscr G(\mathbf L)\) is a formal group topologically flat over
\(k\), whose \(\mathbf O_k\)-bialgebra is \(U(\mathbf L)\).

\item[\textup{(ii)}]
\(\mathscr G(\mathbf L)\) is infinitesimal.

\item[\textup{\SGAInlineRefTarget{sga3:VIIB:2.6.2:prop:item:iii}(iii)}]
For every object \(C\) of \(\operatorname{Alf}/k\),
\(\operatorname{Lie}(\mathscr G(\mathbf L))(C)\) is identified with
the set
\[
  \operatorname{Prim}U(\mathbf L(C))
  =
  \left\{
    x\in U(\mathbf L(C))
    \,\middle|\,
    \varepsilon(x)=0,\quad
    \Delta(x)=x\otimes1+1\otimes x
  \right\}
\]
of primitive elements of \(U(\mathbf L(C))\). In particular, there
is a natural morphism of Lie algebras over \(\mathbf O_k\),
\[
  \tau_{\mathbf L}:
  \mathbf L\longrightarrow
  \operatorname{Lie}\bigl(\mathscr G(\mathbf L)\bigr).
\]
\end{enumerate}}

\textit{Proof.}
The hypothesis that \(\mathbf L\) is flat over \(\mathbf O_k\)
means that, for every morphism \(C\to D\) in
\(\operatorname{Alf}/k\),
\[
  \mathbf L(D)=\mathbf L(C)\otimes_C D,
\]
and that, for every local component \(C'\) of \(C\), the
\(C'\)-module \(\mathbf L(C')\) is free. The first condition,
together with the universal properties of tensor products and
enveloping algebras, gives
\[
  U\bigl(\mathbf L(D)\bigr)
  =
  U\bigl(\mathbf L(C)\bigr)\otimes_C D.
\]
The second condition and the Poincaré--Birkhoff--Witt theorem
\textup{(cf.\ Bourbaki, \emph{Groupes et algèbres de Lie},
I, 2.7)} show that \(U(\mathbf L(C'))\) is a free \(C'\)-module.
Consequently, the bialgebra \(U(\mathbf L)\) is flat over
\(\mathbf O_k\), proving \textup{\SGARefLink{sga3:VIIB:2.6.2:prop:item:i}{(i)}}.

To prove \textup{\SGARefLink{sga3:vii:paragraph:5-1:equation:ii}{(ii)}} and \textup{\SGARefLink{sga3:VIIB:2.6.2:prop:item:iii}{(iii)}}, we may suppose that
\(k\) is Artinian. Put
\[
  L=\mathbf L(k),\qquad
  U=U(L),\qquad
  U_0=k\cdot1_U,
\]
and let \(U^+\) be the two-sided ideal of \(U\) generated by the
image of \(L\). For every \(n\geq1\), put
\[
  U_n
  =
  \left\{
    u\in U
    \,\middle|\,
    \Delta_U(u)-u\otimes1\in U_{n-1}\otimes U^+
  \right\}.
\]
By \SGARefLink{sga3:VIIB:1.3.6}{1.3.6}, it is enough to show that \(U\) is the union of the
\(U_n\). Identify \(L\) with its canonical image in \(U\). Then
\(L\subset U_1\). If \(x_1,\ldots,x_n\in L\), then
\[
  \Delta_U(x_1\cdots x_n)
  =
  (x_1\otimes1+1\otimes x_1)\cdots
  (x_n\otimes1+1\otimes x_n).
\]
Induction on \(n\) therefore shows that
\(x_1\cdots x_n\in U_n\), and hence
\[
  U=\bigcup_n U_n.
\]
This proves \textup{\SGARefLink{sga3:vii:paragraph:5-1:equation:ii}{(ii)}}.

Finally, let
\[
  D=k[d]/(d^2)
\]
be the algebra of dual numbers over \(k\). By hypothesis,
\[
  \mathbf L(D)\simeq L\otimes_kD,
  \qquad
  U(\mathbf L(D))\simeq U\otimes_kD.
\]
It follows that
\(\operatorname{Lie}(\mathscr G(\mathbf L))(k)\) is identified with
the set of elements
\[
  z=1+xd\in U\oplus Ud
  \qquad (x\in U)
\]
such that \(\varepsilon(z)=1\) and \(\Delta(z)=z\otimes z\).
These conditions are equivalent to
\[
  \varepsilon(x)=0,
  \qquad
  \Delta(x)=x\otimes1+1\otimes x,
\]
that is, to \(x\in\operatorname{Prim}U\). In particular,
\[
  \tau_{\mathbf L}:x\longmapsto1+dx
\]
is a morphism of Lie algebras over \(\mathbf O_k\) from
\(\mathbf L\) to
\(\operatorname{Lie}(\mathscr G(\mathbf L))\).

